\subsection{Simulation}

The simulation phase is performed only for the LQR controller with integral action, since this is the
architecture selected for the real implementation. Both the linear and nonlinear Simulink models are
used to check the position response, the current request and the possible activation of current
saturation.

The nominal simulation tuning is
\begin{equation}
    Q_{aug,sim}
    =
    \mathrm{diag}
    \left(
        5000,\;
        5\times10^{-1},\;
        10^{-9},\;
        9\times10^{4}
    \right),
    \qquad
    R_{aug,sim} = 5\times10^{-4}.
\end{equation}

This set of parameters gives a satisfactory response in simulation, where the model is exact and the
measurements are ideal. However, the same tuning is not directly suitable for the real plant: measurement
noise, unmodelled nonlinearities, current saturation and sensor calibration errors require a less
aggressive experimental tuning. Conversely, the parameters selected for the real setup are not the best
choice in simulation, because they are intentionally adapted to the limitations of the physical system.

% \begin{figure}[H]
%     \centering
%     \includegraphics[width=0.85\textwidth]{lqi_step_simulation.pdf}
%     \caption{Simulation response of the LQR controller with integral action.}
%     \label{fig:lqi_step_simulation}
% \end{figure}

\subsubsection{Stability margins}

For the LQR design with integral action, the open-loop system used to compute margins is
\begin{equation}
    L_{lqi}(s) = K_{aug}\left(sI - A_{aug}\right)^{-1} B_{aug}.
\end{equation}
The gain margin, phase margin, and crossover frequencies are reported here based on the values returned
by \texttt{margin} in \texttt{SS\_model\_LQ.m}.
